\begin{figure}[htbp]
	\centering
	\begin{tikzpicture}[x=1.55cm, y=1.55cm,
		trave/.style  ={line width=2.0pt, line cap=round, black},
		incog/.style  ={-{Stealth[length=2.2mm,width=1.3mm]},
			line width=0.9pt, black},
		der/.style    ={-{Stealth[length=1.8mm,width=1.0mm]},
			line width=0.45pt, gray!70},
		terra/.style  ={line width=0.7pt},
		nodoV/.style  ={circle, fill=white, draw=black,
			inner sep=0pt, minimum size=4.2pt, line width=0.6pt},
		nodoI/.style  ={circle, fill=white, draw=black,
			inner sep=0pt, minimum size=6.5pt, line width=0.9pt},
		et/.style     ={font=\footnotesize},
		etd/.style    ={font=\scriptsize, gray!70}]

		\coordinate (A) at (0,0);  \coordinate (B) at (1,-1);
		\coordinate (C) at (1,1);  \coordinate (D) at (1,0);
		\coordinate (E) at (3,0);  \coordinate (F) at (5,0);
		\coordinate (G) at (5,1);  \coordinate (H) at (5,-1);
		\coordinate (L) at (6,0);

		%% ── i corpi, staccati lungo il proprio asse ─────────────
		\draw[trave] (0.297,0.297)--(0.703,0.703);   % ritenuta AC
		\draw[trave] (5.703,0.297)--(5.297,0.703);   % ritenuta LG
		\draw[trave] (1,-0.58)--(1,0.58);            % antenna BC
		\draw[trave] (5,-0.58)--(5,0.58);            % antenna HG
		\draw[trave] (1.376,0.812)--(2.624,0.188);   % strallo CE
		\draw[trave] (4.624,0.812)--(3.376,0.188);   % strallo GE
		\draw[trave] (1.50,0)--(2.58,0);             % impalcato DE
		\draw[trave] (4.50,0)--(3.42,0);             % impalcato FE

		%% ── cerniere interne in D e in F (non isolate) ──────────
		\draw[black, fill=white, line width=0.9pt]
			($(D)+(0,0.07)$) arc[start angle=90,end angle=-90,radius=0.07]--cycle;
		\draw[black, fill=white, line width=0.9pt]
			($(F)+(0,-0.07)$) arc[start angle=-90,end angle=-270,radius=0.07]--cycle;

		%% ── vincoli esterni ─────────────────────────────────────
		\foreach \P in {A,L}{%
			\draw[terra] (\P)--($(\P)+(-0.11,-0.19)$)--($(\P)+(0.11,-0.19)$)--cycle;
			\draw[terra] ($(\P)+(-0.17,-0.19)$)--($(\P)+(0.17,-0.19)$);
			\foreach \t in {-0.15,-0.075,0,0.075,0.15}
				\draw[terra, line width=0.4pt] ($(\P)+(\t,-0.19)$)--++(-0.05,-0.07);}
		\foreach \P in {B,H}{%
			\draw[terra] (\P)--($(\P)+(-0.10,-0.18)$)--($(\P)+(0.10,-0.18)$)--cycle;
			\draw[terra, fill=white, line width=0.5pt]
				($(\P)+(-0.055,-0.215)$) circle(0.035);
			\draw[terra, fill=white, line width=0.5pt]
				($(\P)+( 0.055,-0.215)$) circle(0.035);
			\draw[terra] ($(\P)+(-0.16,-0.25)$)--($(\P)+(0.16,-0.25)$);
			\foreach \t in {-0.14,-0.07,0,0.07,0.14}
				\draw[terra, line width=0.4pt] ($(\P)+(\t,-0.25)$)--++(-0.05,-0.07);}

		%% ── estremi dei corpi e vincoli esterni ─────────────────
		\foreach \n in {(0.297,0.297),(0.703,0.703),(5.703,0.297),(5.297,0.703),
			(1,-0.58),(1,0.58),(5,-0.58),(5,0.58),
			(1.376,0.812),(2.624,0.188),(4.624,0.812),(3.376,0.188),
			(1.50,0),(2.58,0),(4.50,0),(3.42,0),(0,0),(6,0),(1,-1),(5,-1)}
			\node[nodoV] at \n {};

		%% ── i tre nodi isolati come corpi liberi ────────────────
		\node[nodoI] at (C) {};
		\node[nodoI] at (E) {};
		\node[nodoI] at (G) {};

		%% ── le sei incognite primarie ───────────────────────────
		\draw[incog] (0.297,0.297)--(0.085,0.085);
		\draw[incog] (0.703,0.703)--(0.915,0.915);
		\node[et] at (0.30,0.64) {$X_{AC}$};
		\draw[incog] (5.703,0.297)--(5.915,0.085);
		\draw[incog] (5.297,0.703)--(5.085,0.915);
		\node[et] at (5.70,0.64) {$X_{LG}$};
		\draw[incog] (2.624,0.188)--(2.892,0.054);
		\draw[incog] (1.376,0.812)--(1.108,0.946);
		\node[et] at (2.02,0.66) {$X_{CE}$};
		\draw[incog] (3.376,0.188)--(3.108,0.054);
		\draw[incog] (4.624,0.812)--(4.892,0.946);
		\node[et] at (3.98,0.66) {$X_{GE}$};
		\draw[incog] (1,-0.90)--(1,-0.64);
		\node[et, left]  at (0.96,-0.80) {$Y_B$};
		\draw[incog] (5,-0.90)--(5,-0.64);
		\node[et, right] at (5.04,-0.80) {$Y_H$};

		%% ── forze che ne discendono, alle interfacce ────────────
		\draw[der] (1,0)--(0.72,0);         \node[etd, above] at (0.80,0.02){$X_D$};
		\draw[der] (1,0)--(1,-0.28);        \node[etd, left]  at (0.97,-0.20){$Y_D$};
		\draw[der] (1.50,0)--(1.78,0);      \node[etd, below] at (1.72,-0.03){$X_D$};
		\draw[der] (1.50,0)--(1.50,0.28);   \node[etd, right] at (1.56,0.22){$Y_D$};
		\draw[der] (5,0)--(5.28,0);         \node[etd, above] at (5.20,0.02){$X_F$};
		\draw[der] (5,0)--(5,-0.28);        \node[etd, right] at (5.03,-0.20){$Y_F$};
		\draw[der] (4.50,0)--(4.22,0);      \node[etd, below] at (4.28,-0.03){$X_F$};
		\draw[der] (4.50,0)--(4.50,0.28);   \node[etd, left]  at (4.44,0.22){$Y_F$};
		\draw[der] (1,0.58)--(1.28,0.58);   \node[etd, below] at (1.26,0.55){$X_C$};
		\draw[der] (1,0.58)--(1,0.86);      \node[etd, left]  at (0.94,0.70){$Y_C$};
		\draw[der] (5,0.58)--(4.72,0.58);   \node[etd, below] at (4.74,0.55){$X_G$};
		\draw[der] (5,0.58)--(5,0.86);      \node[etd, right] at (5.06,0.70){$Y_G$};
		\draw[der] (2.58,0)--(2.86,0);      \node[etd, below] at (2.80,-0.03){$X_{ED}$};
		\draw[der] (2.58,0)--(2.58,-0.28);  \node[etd, left]  at (2.55,-0.18){$Y_{ED}$};
		\draw[der] (3.42,0)--(3.14,0);      \node[etd, below] at (3.20,-0.03){$X_{EF}$};
		\draw[der] (3.42,0)--(3.42,-0.28);  \node[etd, right] at (3.45,-0.18){$Y_{EF}$};

		%% ── lettere ─────────────────────────────────────────────
		\node[et, above left=1pt]  at (A) {$A$};
		\node[et, above right=2pt] at (B) {$B$};
		\node[et, above left=1pt]  at (C) {$C$};
		\node[et] at (1.16,0.16) {$D$};
		\node[et, above=8pt]       at (E) {$E$};
		\node[et] at (4.84,0.16) {$F$};
		\node[et, above right=1pt] at (G) {$G$};
		\node[et, above left=2pt]  at (H) {$H$};
		\node[et, above right=1pt] at (L) {$L$};

	\end{tikzpicture}
	\caption{La partizione adottata. Ogni corpo è staccato
		lungo il proprio asse, così che i vuoti si aprano dove
		sono i collegamenti e la maglia resti riconoscibile; i
		cerchietti maggiorati in $C$, in $E$ e in $G$ sono i tre
		nodi isolati come corpi liberi autonomi, mentre in $D$ e
		in $F$ la connessione resta un vincolo interno. Le frecce
		rappresentano forze scambiate ai nodi, non
		caratteristiche della sollecitazione. A tratto pieno le
		sei incognite primarie: una per ciascuna biella, che ne
		trasmette il valore inalterato da un estremo all'altro, e
		le due reazioni dei carrelli. A tratto sottile le forze
		alle interfacce fra antenne e impalcati, che pure
		esistono ma discendono dalle prime. I carichi applicati
		sono quelli della \FIGref{fig:strallato_schema}.}
	\label{fig:strallato_esploso}
\end{figure}
